\chapter{Допуск общего носителя фермионной тахионной восприимчивости}
\label{ch:v10-hopf-reservoir-m4-fermion-susceptibility-carrier}

\section{Условный оператор на \texorpdfstring{$M_4$}{M4}-носителе}

Дополнительного линейного пространства не требуется. На
\begin{equation}
 \mathcal H_F=\mathbb C^2_{\rm res}\oplus\mathbb C^2_{\rm path},
 \qquad \dim_{\mathbb C}\mathcal H_F=4
\end{equation}
введём секторную градуировку и cross-инволюцию
\begin{equation}
 R=\begin{pmatrix}I_2&0\\0&-I_2\end{pmatrix},\qquad
 X=\begin{pmatrix}0&I_2\\I_2&0\end{pmatrix}.
\end{equation}
Они образуют точную клиффордову пару:
\begin{equation}
 R^2=X^2=I_4,\qquad RX+XR=0.
\end{equation}
Поэтому семейство
\begin{equation}
 D(x)=R+xX
\end{equation}
удовлетворяет
\begin{equation}
 D(x)^2=(1+x^2)I_4,qquad
 \det D(x)=(1+x^2)^2.
\end{equation}

\section{Детерминант действительно создаёт конденсацию}

Фермионный вклад равен
\begin{equation}
 \Gamma_F(x)=-\log\det D(x)=-2\log(1+x^2),
 \qquad \Gamma_F''(0)=-4.
\end{equation}
Вместе с единичной положительной incidence-жёсткостью получается
\begin{equation}
 \Gamma_{\rm tot}(x)=x^2-2\log(1+x^2),
 \qquad \Gamma_{\rm tot}''(0)=-2.
\end{equation}
Этот функционал ограничен снизу и имеет точные стационарные точки
\begin{equation}
 x\in\{-1,0,1\},\qquad
 \Gamma_{\rm tot}''=(-1,0,1)=(2,-2,2).
\end{equation}
Следовательно, $x=\pm1$ --- два устойчивых минимума со значением
$1-\log4<0$.

Для комплексного фермионного гауссиана антисимметричная Berezin-матрица
на восьми нечётных координатах имеет ранг $8$ и определитель $16$.
Таким образом, условная determinant-архитектура полностью согласована.

\section{Граница унаследованности}

Текущий блок-диагональный родитель содержит носитель и градуировку $R$, но
его проекция cross-инволюции равна
\begin{equation}
 P_{\rm res}XP_{\rm res}+P_{\rm path}XP_{\rm path}=0.
\end{equation}
Кроме того, обычное комплексное пространство не получает автоматически
нечётную статистику, Berezin-меру и нормировку coupling, при которой
фермионная кривизна превосходит incidence-жёсткость. Поэтому условная
конденсация не является физическим выводом.

\begin{theorem}[Минимальный условный фермионный носитель]
Четырёхмерный резервуарно-хопфовский носитель допускает Dirac-family
$D(x)=R+xX$, чей фермионный логарифм определителя создаёт ограниченный
снизу потенциал с двумя ненулевыми устойчивыми минимумами.
\end{theorem}

\begin{nogocriterion}[Комплексный носитель не назначает Grassmann-статистику]
Алгебраическая возможность записать $D(x)$ не выводит cross-билинейный
оператор, parity shift, Berezin-меру или его coupling. Пока эти данные не
следуют из родителя, determinant-механизм остаётся условным.
\end{nogocriterion}

\begin{protocol}\begin{enumerate}
\item минимальная комплексная размерность: \textbf{$4$};
\item условная архитектура: \textbf{$14/14$};
\item фермионная кривизна: \textbf{$-4$};
\item полная кривизна при нуле: \textbf{$-2$};
\item устойчивые минимумы: \textbf{$x=\pm1$};
\item унаследованные ингредиенты: \textbf{$2/6$};
\item строгое физическое происхождение: \textbf{$0/4$};
\item следующий гейт: \textbf{аудит cross-билинейности и нечётной статистики}.
\end{enumerate}\end{protocol}

Машинный сертификат сохранён в
\path{s2t_v10_cell_birth_four_volume_induced_newton_breathing_anomaly_hopf_reservoir_m4_fermionic_determinant_tachyonic_susceptibility_common_carrier_admission_gate_results.json}.